\chapter{Gudang \& Pengadaan (Inventory \& Procurement)}

Modul ini bertanggung jawab untuk menjaga keseimbangan antara keluar masuknya barang, memastikan tidak ada \textit{overselling} (stok negatif), serta mencatat pengadaan dari \textit{Supplier}.

\section{Procurement \& Purchase Order (Pengadaan)}

Untuk menambah stok barang yang kosong, Admin atau staf terkait membuat rencana pengadaan.

\subsection{Skenario Pembelian Barang dari Supplier}
\begin{enumerate}
    \item Buka menu \textbf{Purchase Orders}.
    \item Klik \textbf{New Purchase Order}.
    \item Pilih \textbf{Supplier} (misal: "PT Supplier Bahan Baku").
    \item Tentukan produk dan jumlah (Qty) yang hendak dibeli.
    \item Simpan form tersebut. Status awal adalah \textit{Pending}.
    \item Setelah barang tiba di lokasi fisik, pengguna dengan peran \textbf{Gudang} dapat mengklik \textbf{Terima Barang} untuk mencatat mutasi stok masuk (\textit{Stock In}).
\end{enumerate}

\section{Stock Movements (Mutasi Stok)}

Sistem mencatat setiap aliran pergerakan barang, baik penambahan stok dari Pembelian maupun pengurangan stok dari Pengiriman.

\subsection{Manajemen Stok Presisi (Atomic Lock)}
\textbf{Perhatian Keamanan}: Sistem ini telah dilengkapi \textit{Database Atomic Lock} (\texttt{lockForUpdate}).
Jika ada dua kasir atau dua staf gudang yang memproses pengiriman produk secara bersamaan dalam milidetik yang sama, sistem tidak akan mengalami cacat logika. Sistem akan me-\textit{lock} tabel stok, menyelesaikan transaksi pertama, lalu menghitung saldo terbaru untuk transaksi kedua.

\begin{figure}[H]
    \centering
    \includegraphics[width=0.9\textwidth]{placeholder_stock_movement.png}
    \caption{Daftar Riwayat Pergerakan Stok (In/Out)}
    \label{figure:stock_movement}
\end{figure}
